\documentclass[10pt,a4paper]{amsart}
\usepackage[margin=19mm]{geometry}
\usepackage{amsmath,amssymb,mathtools}
\usepackage[scr=rsfs,cal=euler]{mathalfa}
\usepackage{xfrac}
\usepackage[unicode,hidelinks,bookmarks=false]{hyperref}
\newcommand{\ie}{i.e.\ }
\newcommand{\cf}{cf.\ }
\newcommand{\resp}{respectively\ }
\newcommand{\Subsection}{\subsection}
\providecommand{\nobreakdash}{\nobreak\hbox{-}}
\setlength{\emergencystretch}{2em}
\multlinegap0pt
\newtheorem{lem}{Lemma}
\newtheorem{thm}{Theorem}
\def\A{\textrm{A}}
\def\C{\textrm{C}}
\def\V{\mathcal V}
\def\X{\mathcal X}
\def\cl{\textrm{cl}}
\def\cogen{cogen}
\def\cogenz{\cogen_0}
\def\eq{eq}
\def\gen{gen}
\def\genz{\gen_0}
\title{Completeness of Relational Algebra via Cylindric Algebra}
\author{Jan La\v{s}tovi\v{c}ka}
\date{Source: arXiv:2603.15099v1 (CC BY 4.0)}
\begin{document}
\maketitle

\noindent\emph{Source and licence.} Completeness of Relational Algebra via Cylindric Algebra. Version arXiv:2603.15099v1 (CC BY 4.0), distributed by arXiv under the Creative Commons Attribution 4.0 International licence, http://creativecommons.org/licenses/by/4.0/, as recorded in the arXiv licence metadata for this exact version and shown on the abstract page https://arxiv.org/abs/2603.15099v1. Full licence notice: \texttt{licenses/arxiv-2603.15099-CC-BY-4.0.txt}.\par\smallskip

\noindent\textit{Editorial scope.} Counted unit: the article's thm thm (source0.tex lines 571--577). The article's complete original proof of it is delivered with it (source0.tex lines 578--638, one \texttt{proof} environment of 61 lines). No mathematics is written, repaired or re-derived: the statement and the proof are the article's own lines, in the article's own notation, and no retained line is substituted or re-flowed.  Retained uncounted as the source-local context the proof needs: lines 533--538; lines 544--549; lines 554--559; lines 565--568.  Nothing was omitted from any retained span, so the package carries no omission record.  No retained line contains a citation, so the wrapper carries no bibliography and no bibliography entry.  No retained line contains a figure environment, an image-inclusion command or drawing code, so no image is delivered and the package declares no asset dependency.  Editorial formatting only, in addition: an AMS \texttt{amsart} preamble that restates the article's own declarations and macros used by the retained lines, namely the environments \texttt{lem}, \texttt{thm} on separate counters, together with the standard packages those lines need.  The retained environments print in this document as Lemma 1, Theorem 1 in order of appearance, whereas the article prints the same items with the numbering of its own full document.  The complete native source of the article as distributed in the arXiv e-print for this exact version is bundled unchanged as \texttt{sources/196458/source0.tex}.
\begin{lem}\label{lem:cap_project}
Let $X_1,X_2\subseteq\X$ and $V_1,V_2\subseteq\V$. Then it holds
\begin{align*}
\epsilon(\pi_{X_1\cup X_2}(V_1\cap V_2))\subseteq\epsilon(\pi_{X_1}(V_1))\cap\epsilon(\pi_{X_2}(V_2)).
\end{align*}
\end{lem}

\begin{lem}\label{lem:cap_dom}
Let $N\subseteq M$, $X_1,X_2\subseteq\X$. Then it holds
\begin{align*}
\epsilon(N^{X_1\cup X_2})=\epsilon(N^{X_1})\cap\epsilon(N^{X_2}). 
\end{align*}
\end{lem}

\begin{lem}\label{lem:eq_dom}
Let $N\subseteq M$, $X\subseteq\X$, and $\varphi$ by a formula such that $\pi_X(||\varphi||)\subseteq N^X$. Then it holds
\begin{align*}
\pi_{\cl_{\eq(\varphi)}(X)}(||\varphi||)\subseteq N^{\cl_{\eq(\varphi)}(X)}. 
\end{align*}
\end{lem}

\begin{align}
&\pi_Y(\cl_X(V))=\pi_Y(V),&\textrm{ provided that } X\cap Y=\emptyset, \label{eq:pi_C_disjoint}\\
&\pi_Y(\overline{\cl_X(V)})\subseteq\pi_Y(\overline{V}).& \label{eq:pi_overline_C}
\end{align}

\begin{thm}
For every formula $\varphi$, it holds
\begin{align*}
\pi_{\genz(\varphi)}(||\varphi||)&\subseteq \A(M)^{\genz(\varphi)},\\
\pi_{\cogenz(\varphi)}(\overline{||\varphi||})&\subseteq \A(M)^{\cogenz(\varphi)}. 
\end{align*}
\end{thm}

\begin{proof}
We prove the assertion by structural induction on formulas. Clearly, it holds for any atomic formula. Suppose that the assertion holds for formulas $\varphi$ and $\psi$. We have
\begin{align*}
\pi_{\genz(\neg\varphi)}(||\neg\varphi||)=\pi_{\cogenz(\varphi)}(\overline{||\varphi||})\subseteq\A(M)^{\cogenz(\varphi)}=\A(M)^{\genz(\neg\varphi)}
\end{align*}
and 
\begin{align*}
\pi_{\cogenz(\neg\varphi)}(\overline{||\neg\varphi||})&=\pi_{\genz(\varphi)}(\overline{\overline{||\varphi||}})=\pi_{\genz(\varphi)}(||\varphi|||)\\
&\subseteq\A(M)^{\genz(\varphi)}=\A(M)^{\cogenz(\neg\varphi)}.
\end{align*}
Therefore, the assertion holds for negated formulas.
From the assumptions $\pi_{\genz(\varphi)}(||\varphi||)\subseteq \A(M)^{\genz(\varphi)}$ and $\pi_{\genz(\psi)}(||\psi||)\subseteq \A(M)^{\genz(\psi)}$, it follows that 
\begin{align*}
\epsilon(\pi_{\genz(\varphi)}(||\varphi||))\cap\epsilon(\pi_{\genz(\psi)}(||\psi||))\subseteq\epsilon(\A(M)^{\genz(\varphi)})\cap\epsilon(\A(M)^{\genz(\psi)}).
\end{align*}
From Lemmas \ref{lem:cap_project} and \ref{lem:cap_dom}, it follows that
\begin{align*}
\pi_{\genz(\varphi)\cup\genz(\psi)}(||\varphi\wedge\psi||)\subseteq\A(M)^{\genz(\varphi)\cup\genz(\psi)}. 
\end{align*}
From Lemma \ref{lem:eq_dom}, we obtain
\begin{align*}
\pi_{\genz(\varphi\wedge\psi)}(||\varphi\wedge\psi||)\subseteq\A(M)^{\genz(\varphi\wedge\psi)}. 
\end{align*}
From the assumptions 
$$\pi_{\cogenz(\varphi)}(\overline{||\varphi||})\subseteq \A(M)^{\cogenz(\varphi)}$$ 
and 
$$\pi_{\cogenz(\psi)}(\overline{||\psi||})\subseteq \A(M)^{\cogenz(\psi)}$$
together with the equality $\cogenz(\varphi\wedge\psi)=\cogenz(\varphi)\cap\cogenz(\psi)$, we obtain 
\begin{align*}
\pi_{\cogenz(\varphi\wedge\psi)}(\overline{||\varphi||})&\subseteq \A(M)^{\cogenz(\varphi\wedge\psi)},\\
\pi_{\cogenz(\varphi\wedge\psi)}(\overline{||\psi||})&\subseteq \A(M)^{\cogenz(\varphi\wedge\psi)}. 
\end{align*}
It follows that
\begin{align*}
\pi_{\cogenz(\varphi\wedge\psi)}(\overline{||\varphi\wedge\psi||})&= \pi_{\cogenz(\varphi\wedge\psi)}(\overline{||\varphi||}\cup\overline{||\psi||})\\
&=\pi_{\cogenz(\varphi\wedge\psi)}(\overline{||\varphi||})\cup\pi_{\cogenz(\varphi\wedge\psi)}(\overline{||\psi||})\\
&\subseteq\A(M)^{\cogenz(\varphi\wedge\psi)}.
\end{align*}
We have shown that the assertion also holds for conjoined formulas.

From the assumption $\pi_{\genz(\varphi)}(||\varphi||)\subseteq \A(M)^{\genz(\varphi)}$, we obtain
\begin{align*}
\pi_{\genz(\varphi)\setminus\{x\}}(||\varphi||)\subseteq \A(M)^{\genz(\varphi)\setminus\{x\}}. 
\end{align*}
Using the equality \eqref{eq:pi_C_disjoint}, it follows that
\begin{align*}
\pi_{\genz((\exists x)\varphi)}(||(\exists x)\varphi||)&=\pi_{\genz(\varphi)\setminus\{x\}}(\C_x(||\varphi||))=\pi_{\genz(\varphi)\setminus\{x\}}(||\varphi||)\\
&\subseteq\A(M)^{\genz(\varphi)\setminus\{x\}}=\A(M)^{\genz((\exists x)\varphi)}.
\end{align*}

From the assumption $\pi_{\cogenz(\varphi)}(\overline{||\varphi||})\subseteq \A(M)^{\cogenz(\varphi)}$, we obtain
\begin{align*}
\pi_{\cogenz(\varphi)\setminus\{x\}}(\overline{||\varphi||})\subseteq \A(M)^{\cogenz(\varphi)\setminus\{x\}}. 
\end{align*}
Using the equality \eqref{eq:pi_overline_C}, it follows that
\begin{align*}
\pi_{\cogenz((\exists x)\varphi)}(\overline{||(\exists x)\varphi||})&=\pi_{\cogenz(\varphi)\setminus\{x\}}(\overline{\C_x(||\varphi||)})\subseteq\pi_{\cogenz(\varphi)\setminus\{x\}}(||\varphi||)\\
&\subseteq\A(M)^{\cogenz(\varphi)\setminus\{x\}}=\A(M)^{\cogenz((\exists x)\varphi)}.
\end{align*}
We have shown that the assertion also holds for existentially quantified formulas.
\end{proof}

\end{document}
